\documentclass[11pt]{article}

\newcommand{\cnum}{Math 245B}
\newcommand{\ced}{Winter 2026}
\newcommand{\ctitle}[4]{\title{\vspace{-0.5in}\cnum, \ced\\Week #1: #2}\author{\vspace{-0.35in}\\#3, #4}}
\usepackage{enumitem}
\usepackage[usenames,dvipsnames,svgnames,table,hyperref]{xcolor}
\usepackage{amsmath, amsfonts}
\usepackage{graphicx} % Required for inserting images
\usepackage{float}
\usepackage{listings}
\usepackage{xcolor}
\usepackage{hyperref}
\usepackage{comment}

\renewcommand*{\theenumi}{\alph{enumi}}
\renewcommand*\labelenumi{(\theenumi)}
\renewcommand*{\theenumii}{\roman{enumii}}
\renewcommand*\labelenumii{\theenumii.}

\definecolor{codegreen}{rgb}{0,0.6,0}
\definecolor{codegray}{rgb}{0.5,0.5,0.5}
\definecolor{codepurple}{rgb}{0.58,0,0.82}
\definecolor{backcolour}{rgb}{0.95,0.95,0.92}
\definecolor{header}{HTML}{205459}
\definecolor{solution}{HTML}{204d52}

\newcommand{\solution}[1]{{{\textcolor{header}{Solution:} \textcolor{solution}{#1}}}}
\setlist[enumerate]{label=(\alph*)}

\lstdefinestyle{mystyle}{
    backgroundcolor=\color{backcolour},
    commentstyle=\color{codegreen},
    keywordstyle=\color{magenta},
    numberstyle=\tiny\color{codegray},
    stringstyle=\color{codepurple},
    basicstyle=\ttfamily\footnotesize,
    breakatwhitespace=false,
    breaklines=true,
    captionpos=b,
    keepspaces=true,
    numbers=left,
    numbersep=5pt,
    showspaces=false,
    showstringspaces=false,
    showtabs=false,
    tabsize=2
}


\begin{document}
\ctitle{\#2}{}{Asher Christian}{006-150-286}
\date{}
\maketitle
\section{Corollary}
\emph{
    If $f(z) = \sum_{n=0}^{\infty}\alpha_n z^{n}$ converges in $\{|z| < R\}, R > 0$,then
    $f$ is holomorphic in $\{|z| < R\}$ and  $f$ is infinitely differentiable in its radius of convergence and
    \[
    f^{(k)}(z) = \sum_{n=k}^{\infty}n(n-1)\dots(n-k+1)\alpha_nz^{n-k}
    \] 
    and
    \[
    \alpha_n = \frac{f^{(n)}(0)}{n!}
    \]  
}\\
\emph{
    Definition:\\
    $f(z)$ is analytic on the domain  $\Omega$ if $\forall z_0 \in \Omega$ $f$ has power series expansion
    \[
    \sum_{}^{}\alpha_k(z_0)(z-z_0)^{k}
    \] 
    converging in $\{|z-z_0| < \delta_{z_0}\}$ $\delta_{z_0} > 0$
}
so $f$ analytic in $\Omega$ implies  $f$ is holomorphic in $\Omega$.
Ex: 
 \[
f(x) =
\begin{cases}
    x^{\frac{3}{2}} & x \ge 0 \\
    0 & x < 0
\end{cases}
\] 
\[
f'(x) = 
\begin{cases}
    \frac{3}{2}x^{\frac{1}{2}} & x \ge 0\\
    0 & x < 0
\end{cases}
\] 

\section{Theorem}
For $\omega \in \mathbb{C}\setminus \{0\}$ and $D_{\omega} = \{z \in \mathbb{C}: |z-w| < |w|\}$
Then there exists $g : D_\omega \rightarrow \mathbb{C}$ such that $e^{g(z)} = z$  for all $z \in D_\omega$

\section{Tuesday}
\emph{
    Def: A parametized curve is a continuous function $z : [a,b] \rightarrow \mathbb{C}$ $-\infty < a < b < \infty$.
    It is smooth if $z'(t)$ exists and is continuous and $z'(t) \ne 0$ for all  $a < t < b$ and
     \[
    z'(a) = \lim_{t\downarrow 0} \frac{z(a+t) - (a)}{t}
    \] 
    $z(t)$ is piecewise smooth if  $z(t)$ is continuous and there exists  $a = a_0 < a_1 < \dots < a_m = b$ and
    $z'(t)$ is continuous on  $(a_{j-1},a_j)$  $1 \le j \le m$. Left and right derivatives exist at  $a_j$ for all  $j$
}
\\
\emph{
    Def: $z : [a,b] \rightarrow \mathbb{C}$ and $\tilde z: [c,d] \rightarrow \mathbb{C}$ are equivalent
    if there exists a map 
    \[
        t(s) : [c,d] \rightarrow^{\text{onto}} [a,b]
    \] 
    $t'(s) > 0$ so that  $z(t(s)) = \tilde z(s)$
}
\\
Ex:
\[
    z_n : (0,2\pi] \rightarrow \mathbb{C}
\] 
\[
z(t)_n(t) = e^{int}
\] 

The reverse curve of $z : [a,b] \rightarrow \mathbb{C}$ is $z^{-}(t) : [a,b] \rightarrow \mathbb{C}$ but $z^{-}(t) = z(b+a-t)$

Let
\[
    \gamma = \{z(t) : a \le t \le b\}
\] 
be a parametized line
\[
f : \gamma \rightarrow \mathbb{C}
\] 
be continuous. Define
\[
\int_{\gamma}^{}f(z)dz = \int_{a}^{b}f(z(t))z'(t)dz
\] 

\section{Lemma 4.1}
\emph{
    \begin{enumerate}
        \item 
            \[
            \int_{\gamma}^{}f(z)dz
            \] 
            is independent of the equivalent parametization of $\gamma$.
            $\gamma = z(t)$ $a \le t \le b$
             $\tilde \gamma = \tilde z(s)$ $c \le s \le d$ then
              \[
             \int_{\gamma}^{}fdz = \int_{\tilde \gamma}^{}fdz
             \] 
             \[
                 \frac{d \tilde z}{ds} = \frac{dz}{dt} \frac{dt}{ds}
             \] 
         \item 
             \[
             \int_{- \gamma}^{}f(z)dz = - \int_{\gamma}^{}f(z)dz
             \] 
             \[
             z(t) = z(t(s)) = \tilde z(s)
             \] 
    \end{enumerate}
}

\section{Wednesday}
\section{Theorem 4.3}
$F$ holo in a domain $\Omega$. Assume  $F'(z)$ is continuous in  $\Omega$. Let $\gamma = \{z(t) : a \le t \le b\}$ piecewise
smooth parametized curve.
\[
\int_{\gamma}^{}F'(z)dz = F(z(b)) - F(z(a))
\] 
\section{Corollary}
If $F'$ is continuous and  $\gamma$ is closed i.e. $z(b) = z(a)$ then
 \[
\int_{\gamma}^{}F'(z)dz = 0
\] 
If $F' = 0$ on  $\Omega$ then  $F$ is constant. Ex:  $T$ is a solid triangle and $T \subset \Omega$.
Triangle is oriented counter clockwise. Then
\[
\int_{\partial T}^{}p(z)dz = 0
\] 
for all polynomials $p$

\section{Cauchy integral theorem}
\emph{
    Let $ D$ be a disc.  $\gamma$ piecewise smooth simple closed curve. $\gamma \subset D$. $f$ holomorphic in $D$  implies
     \[
    \int_{\gamma}^{}f(z)dz = 0
    \]
}
$f(z) = \frac{1}{z}$ is holomorphic in $ \mathbb{C} \setminus \{0\}$ but $\int_{ \gamma}^{}f(z)dz \ne 0$ $\int_{0}^{2\pi}e^{-i\theta}e^{i\theta}d\theta = 2\pi$
\\
\emph{
    $f$ holomorphic on domain $\Omega \implies f'$ is continuous in  $\Omega$
}
\\
\emph{
    $f$ holomorphic implies $f$ is analytic (i.e. has powerseries experession at all points)
}
\\
\emph{
    if $\Omega$ is a domain and  $f$ is holomorphic on $\Omega$.  $D = \{z : |z-z_0r < r\} \subset \Omega$ then $f|_{D}$ determins  $f$
}
\\
\emph{
    normal families arguments (compactness)
}

\section{Theorem (goursot)}
\emph{
    Let $T$ be a closed solid triangle contained in a domain  $\Omega$ and let  $f$ be holomorphic on $\Omega$.
    Then
    \[
    \int_{\partial T}^{}f(z)dz = 0
    \] 
}
\\Proof:
Write $T^{0} = T$ $T_{J_1}^{1}, j=1,2,3,4$  are formed by connecting the midpoints of sides of $T^{0}$.
each $T_{J_1}^{(1)}$ gives rise to $T_{j_1,j_2}^{(2)}$
 for $j_2=1,2,3,4$ and so we subdivide each other traingle into four parts via the same technique.
 In general $T_{j_1,\dots,j_N}^{(N)}$ gives rise to $T_{j_1,\dots,j_N,j_{N+1}}^{(n+1)}$.
 THen
 $l(\partial T_{j_1,\dots,j_n}^{(n)} = 2^{-m} l(\partial T^{0})$ orean each 
 $\partial(T_{j_1,\dots,j_n}^{(n)})$ coounter clockwise.
 then
 \[
     \int_{\partial T^{0}}^{}f(z)dz = \sum_{j=1}^{4}\int_{\partial T_{j}^{(1)}}^{}f(z)dz
 \] 

so
\[
    \int_{\partial T_{\vec j}^{(n)}}^{} = \sum_{j_{n+1}}^{4}\int_{\partial T_{\vec j,j_{n+1}}^{(n+1)}}^{}f(z)dz
\] 
Assume that
\[
\left|\int_{\partial T}^{}f(z)dz\right| = \alpha > 0
\] 
then exists $j_1$ such that $\left | \int_{\partial T_{j_1}^{1}}^{}f(z)dz \right| > \frac{\alpha}{4}$.
THen there exists $j_1,j_2,\dots,j_n$ such that
\[
    \left | \int_{\partial T_{j_1,\dots,j_n}^{(n)}}^{}f(z)dz \right| \ge \frac{\alpha}{4^{n}}
\] 
by math 131
 \[
     \bigcap_{n=1}^{\infty}T_{j_1,j_2,\dots,j_n}^{n} =\{z_0\}, z_0 \in \Omega
\] 
$f(z) = f(z_0) + f'(z_0)(z-z_0) + o(|z-z_0|)$ $R(z) = f(z) - p(z)$ with  $p(z) = f(z_0) + f'(z_0)(z_-z_0)$. then
\[
    \sup_{T^{(n)}_{j_1,\dots,j_n}}|R(z)| \le \epsilon 2^{-n}
\] 
since the max width of $T^{(n)}$ is $2^{-n}$
if $n \ge n(\epsilon)$\\
therefore  since the integral of a polynomial is known to be zero around a closed curve
\[
    \left |\int_{\partial T_{j_1,\dots,j_n}^{(n)}}^{}f(z)dz \right| = \left | \int_{\partial T_{j_1,\dots,j_n}^{(n)}}^{} R(z)dz\right | \le \epsilon 2^{-n}l(T)2^{-n}
\] 
so we get
\[
\frac{\alpha}{4^{n}} < \epsilon l(T) \frac{1}{4^{n}}
\] 
but $\alpha$ is fixed and we can pick any $\epsilon$ we want therefore
$\alpha = 0$ so we are done.

\section{Theorem 5.2}
\emph{
    If $f(z)$ is holomorphic in an open disc  $D$. 
    then  $\exists$  $F(z)$ holomorphic in  $D$ with  $F'(z) = f(z)$ in  $D$. NOT TRUE FOR ALL DOMAINS  $\Omega$.
}

\section{Friday Cauchy integral theorem}
\emph{
   $f$ holomorphic in a disc $ D$ and $T$ triangle subset of  $D$ then
    \[
   \int_{\partial T}^{}f(z)dz = 0
   \] 
}\emph{\\
   If $f$ is holomorphic in a disc $D$ then there exists $F$ holomorphic in $D$ with  $F'(z) = f(z)$ for all  $z \in D$.
}\\
Proof: define $F(z) = \int_{ \gamma_z}^{}f(w)dw$ where $\gamma_z = \{tz + (1-t)z_0: 0 \le t \le 1\}$. $F(z)$ is defined for all  $z \in D$. also  $z_0$ is the center of $D$.
Also $\sigma := \{tz + (1-t)(z+h): 0 \le t \le 1\}$ I want to show
\[
F(z+h) - F(z)
\] 
from theorem we have
\[
F(z+h) - F(z) + \int_{\sigma}^{}f(w)dw = 0
\] 
so
\[
    \frac{F(z+h) - F(z)}{h} - f(z) = -\frac{1}{h}\int_{\sigma}^{}f(w)dw - f(z)
\] 
RHS
\[
=\frac{1}{h}\int_{\sigma}^{}(f(z)-f(w))dw
\] 
so
\[
    \frac{1}{|h|}|\int_{\sigma}^{}((f(z)-f(w))| \le \frac{1}{|h|}\sup_{w \in \sigma}|f(w)-f(z)| |h|
\] 
and the last term goes to zero since $f$ is continuous. We use $f$ is holomorphic to show that the integral around the triangle is zero.

\section{Corollary 5.3}
\emph{
    If $f$ is holo on an open disc $D$ and $\gamma$ is a piecewise smooth closed curve in  $D$. THen
    \[
    \int_{\gamma}^{}f(z)dz = 0
    \] 
    since
    \[
    \int_{\gamma}^{}f(z)dz = \int_{\gamma}^{}F'(z)dz = 0
    \] 
}

\section{Extensions of the theorem}
A satr like domain is a domain $D$ such that there exists  $z_0 \in D$ such that for all $z \in D$
 \[
     \{tz + (1-t)z_0 : 0 \le t \le 1\} \subset D
\] 
Theorem $5.2$ is true for all star shaped domain.

\section{Theorem 5.4 Cauchy integral formula}
\emph{
    If $f(z)$ is holomorphic on a domain  $\Omega$ and if  $D$ is a disc with $\overline{D} \subset \Omega$. Then
    \[
        \frac{1}{2\pi i} \int_{\partial D \text{counter clockwise}}^{}\frac{f(w)}{w-z}dw = 
        \begin{cases}
            f(z) & z \in D\\
            0 & z \not\in \overline{D}\\
            \text{undefined} & z \in \partial D
        \end{cases}
    \] 
}\\
Proof: if $z \not\in \overline{ D}$  us corollary $5.3$ for 
 \[
g(w) = \frac{f(w)}{w-z}
\] 
and $\gamma = \partial D$. $g(w)$ is holomorphic on  $D$ since  $w-z = 0 $ only when  $w \not\in D$ so the integral is zero.\\
If  $z \in D$, let $z_0$ be the center of the disc then say $D = \{z : (z-z_0) < R\}$ then there exists $ R' > R$ such that  $\{z : (z-z_0) \le R'\} \subset \Omega$.
Then
\[
    \Omega = \overline{D} \setminus [z,z^{*}], \qquad z^{*} \in \partial \overline{D}
\] 
\[
g(w) = \frac{f(w)}{w-z}
\] 
is holomorphic in $\Omega$ 
Let
\[
\gamma = \Gamma_\epsilon + L' + \gamma_\epsilon' + L
\] 
but
\[
\frac{1}{2\pi i} \int_{\gamma}^{} \frac{f(w)}{w-z}dw = 0
\] 
\[
\frac{1}{2\pi i} \int_{\Gamma_\epsilon}^{}\frac{f(w)}{w-z} dz \rightarrow \frac{1}{2\pi i}\int_{\partial D}^{}\frac{f(w)}{w-z}dw
\] 
\[
\frac{1}{2\pi i} \int_{\gamma_\epsilon}^{}\frac{f(w)}{w-z}dw \rightarrow f(z)
\] 
but
\[
\int_{L}^{}f + \int_{L'}^{}f \rightarrow 0
\] 
since $f$ is continuous


\end{document}
